% Chapter 2: Function Spaces, Duality, and Sobolev Regularity
% Part I of the book.
% STATUS: book-standard prose (Part I).

\chapter{Function Spaces, Duality, and Sobolev Regularity}
\label{ch:part1-function-spaces-duality-and-sobolev-regula}
\label{sec:function_spaces}
% -----------------------------------------------------------------------------

Chapter~\chapref{ch:part1-introduction-the-paradigm-shift-to-infinit} ended by naming the object this book studies: a solution operator $\mathcal{G} : \mathcal{A} \to \mathcal{U}$ that carries an input function---a permeability field, an initial velocity, a boundary forcing---to an output function, and that we intend to learn from data. Naming it is not yet understanding it. The inequality $\|\mathcal{G}(a) - \mathcal{G}_\theta(a)\|_{\mathcal{U}} < \varepsilon$ asserts something about a norm we have not written down, on inputs whose size we have not measured, for a map whose continuity we have not checked. Those decisions are not decorative. Replace $L^2$ by $H^{-1}$ and the residual of a differential equation becomes far easier to make small. Move the input space from $L^\infty$ to $H^s$ with $s > d/2$ and a pointwise sensor reading changes from an unbounded expression into a continuous linear functional, hence into something a learning theory can discuss at all. A family merely bounded in one of these norms may, finally, be precompact in another---what every uniform approximation theorem in this Part quietly demands.

The obstacle is infinite dimensionality. In $\mathbb{R}^d$ the Heine--Borel theorem makes closed, bounded, and compact nearly the same hypothesis, so approximation theory gets its compact domains free whenever the data are bounded. Function spaces forfeit this: Riesz's lemma says the closed unit ball of an infinite-dimensional normed space is never compact. Compactness must therefore be purchased, and the three certificates below---Arzel\`a--Ascoli, Kolmogorov--Riesz, Rellich--Kondrachov---purchase it the same way: a bound measuring smoothness or frequency decay, traded for total boundedness in a weaker norm.

The reader should arrive with Lebesgue measure and the $L^p$ spaces, the elementary theory of normed linear spaces, and Fourier series on the periodic box; duality at the level of Brezis~\cite{brezis2011functional} and Sobolev spaces at the level of Evans~\cite{evans2010partial} help but are rebuilt here as needed, and no distribution theory beyond one definition of a weak derivative is required. Notation follows Table~\ref{tab:notation-master} throughout---$\Omega \subset \mathbb{R}^d$, the multi-index derivative $D^\alpha u$, the torus $\mathbb{T}^d$, the dual pairing $\langle \cdot, \cdot \rangle_{X^*, X}$. The route runs from Banach and Hilbert spaces, through duality and the Gelfand triple, to Sobolev spaces stated with their conditions on $s$, $p$, and $d$ exposed, and finally to Bochner spaces, where time-dependent problems live.

Placement within Part~\ref{part:1} dictates the selection. The Darcy problem of \chapref{ch:part1-partial-differential-equations-as-solution} tests its data against $H^{-1}$ and its solutions in $H^1_0$; the universality theorems of \chapref{ch:part1-operator-approximation-theory-and-universa} quantify over compact subsets of the input space; the spectral convergence rate of \chapref{ch:part1-the-fourier-neural-operator-spectral-theor} is a Fourier-tail estimate of the kind derived below; the Sobolev-regularized loss of \chapref{ch:part1-physics-informed-operators-and-sobolev-reg} leans on the embedding $H^s \hookrightarrow L^\infty$ valid when $s > d/2$. Nothing here is original; the accounting is, and it lets the reader tell, when a result from Part~\ref{part:2} or Part~\ref{part:3} is invoked, which assumption it inherits.

% -----------------------------------------------------------------------------
\section{Banach and Hilbert Spaces}
\index{Banach space}\index{Hilbert space}\index{completeness}\index{Riesz--Fischer theorem}

A normed vector space $(X, \|\cdot\|_X)$ is a \textbf{Banach space} if it is complete for the metric $d(u,v) = \|u - v\|_X$ induced by its norm: every Cauchy sequence converges to an element of $X$. If the norm arises from an inner product $\langle \cdot, \cdot \rangle_X$ through $\|u\|_X = \sqrt{\langle u, u \rangle_X}$, the space is a \textbf{Hilbert space}. The distinction is not terminological. A norm comes from an inner product if and only if it satisfies the parallelogram law, and among the Lebesgue spaces $L^p(\Omega)$ only $p = 2$ passes that test. Orthogonality, orthogonal projections, and the ``best approximation lies in the subspace'' reasoning that makes spectral methods clean are therefore available in $L^2$ and unavailable in $L^1$ or $L^\infty$; conversely, the losses practitioners actually minimize are $L^2$-based precisely because the Hilbert structure makes their gradients tractable.

Completeness is the hypothesis that makes limits of approximating sequences exist inside the space rather than somewhere outside it, and for the spaces of this book it is a theorem rather than an assumption: the Riesz--Fischer theorem states that $L^p(\Omega)$ is complete for every $1 \le p \le \infty$, and the Sobolev spaces introduced below are closed subspaces of products of $L^p$ spaces, hence Banach, hence Hilbert when $p = 2$~\cite{brezis2011functional}. That closedness is what converts a minimizing or approximating sequence into a limit that is again a field in the same space, with its norm bound preserved in the limit.

The single Hilbert-space fact that recurs most often in what follows is that the dual is the space itself.

\begin{theorem}[Riesz Representation Theorem]
\label{thm:ch02-riesz-rep}
Let $H$ be a real Hilbert space and let $H^*$ denote its topological dual. For every bounded linear functional $\ell \in H^*$, there exists a unique element $u_\ell \in H$ such that
\begin{equation}
  \ell(v) = \langle u_\ell, v \rangle_H \quad \forall v \in H, \qquad \text{with} \quad \|\ell\|_{H^*} = \|u_\ell\|_H.
  \label{eq:ch02-riesz}
\end{equation}
\end{theorem}

The identity in \eqref{eq:ch02-riesz} is isometric, not merely bijective: measuring a functional and measuring its representer give the same number. Two consequences organize this book's error analysis. First, in a Hilbert target space there is no meaningful difference between an error field and the error functional it induces, so the residual $\mathcal{G}(a) - \mathcal{G}_\theta(a)$ can be read either as an element of $\mathcal{U}$ or as a bounded linear functional on $\mathcal{U}$ without changing its size. Second, Theorem~\ref{thm:ch02-riesz-rep} supplies the Riesz representer of $\ell \in H^*$ even when $\ell$ was defined by an $L^\infty$-type quantity such as a point evaluation, provided $\ell$ is continuous---which, as \eqref{eq:ch02-eval} below shows, is exactly the question of how smooth the space has to be. Banach spaces have no such identification; there the dual is genuinely a different space, and that asymmetry is the reason negative-order Sobolev spaces exist.

Compactness is where finite-dimensional intuition fails, and the failure is total. Given a proper closed subspace $Y \subsetneq X$ and any $\theta < 1$, Riesz's lemma produces a unit vector whose distance from $Y$ exceeds $\theta$; applying this along an increasing chain of subspaces yields a sequence in the closed unit ball whose pairwise distances are bounded below, hence with no Cauchy subsequence. Boundedness therefore buys nothing by itself~\cite{brezis2011functional}. A set $K \subset X$ is \textbf{compact} when every sequence in $K$ has a subsequence converging to a point of $K$; all the spaces here are metric, so this sequential form is the one in use, and the closed unit ball of $X$ is compact exactly when $X$ is finite-dimensional. The usable substitutes come with extra hypotheses. In $C(\overline{\Omega})$ with the uniform norm, the Arzel\`a--Ascoli theorem characterizes relatively compact sets as those that are uniformly bounded and equicontinuous~\cite{brezis2011functional}. In $L^p(\mathbb{R}^d)$, $1 \le p < \infty$, the Kolmogorov--Riesz criterion says that a family $\mathcal{F}$ is relatively compact if and only if it is bounded, uniformly tight in the sense that
$\sup_{u \in \mathcal{F}} \|u\|_{L^p(|x| > R)} \to 0$ as $R \to \infty$, and equicontinuous in the mean,
$\sup_{u \in \mathcal{F}} \|u(\cdot + h) - u(\cdot)\|_{L^p} \to 0$ as $h \to 0$~\cite{brezis2011functional}.
\index{Arzela--Ascoli theorem}\index{Kolmogorov--Riesz compactness criterion}\index{compact set!in a normed space}\index{precompact set}\index{equicontinuity in the mean}

Read the last condition carefully, because it is the hinge of the whole chapter. A uniform bound $\sup_{u \in \mathcal{F}} \|u\|_{W^{1,p}} \le M$ implies it through the estimate
$\|u(\cdot + h) - u(\cdot)\|_{L^p} \le |h| \|\nabla u\|_{L^p}$,
and more generally a bound in $H^s$ for any $s > 0$ implies it at a fractional rate. Smoothness, not size, is what generates that translation control, but on the whole space it buys no part of the tightness condition, since mass can still escape to infinity. Every compactness statement in this book therefore lives on a bounded domain or on the torus, where tightness is automatic and smoothness alone does the work; on $\mathbb{T}^d$ the statement reads as uniform decay of high-frequency Fourier content, which is the formulation the spectral chapters prefer. This is the precise sense in which compact subsets of function spaces must possess decay of their high-frequency components, and it is why the operator-learning results later in this Part are stated on compact sets of inputs rather than on bounded ones: the universality theorem of Chen and Chen~\cite{chen1995universal} and its neural-operator counterpart in \chapref{ch:part1-operator-approximation-theory-and-universa} both require compactness of the domain set, and on bounded-but-not-compact sets their hypotheses simply are not met. \index{metric entropy}The metric entropy of a Sobolev ball is the quantitative shadow of this fact: the estimate $\log \mathcal{N}(\varepsilon, B_{H^s}, L^2) \sim \varepsilon^{-d/s}$, where $\mathcal{N}$ counts the $L^2$ balls of radius $\varepsilon$ needed to cover the unit ball $B_{H^s}$, replaces the $\varepsilon^{-d}$ covering exponent of $\mathbb{R}^d$, and it governs how many training samples a rate theorem can demand. \chapref{ch:part1-open-theoretical-problems-and-mathematical} reuses it as the estimation term of its risk bound.

% -----------------------------------------------------------------------------
\section{Duality, Dual Pairs, and Gelfand Triples}
\index{topological dual}\index{duality pairing}\index{reflexive space}\index{Banach--Alaoglu theorem}\index{Gelfand triple}

For a Banach space $X$, the \textbf{topological dual} $X^* := \mathcal{L}(X, \mathbb{R})$ is the space of continuous linear functionals on $X$, equipped with the operator norm
\begin{equation}
  \|\ell\|_{X^*} = \sup_{\|v\|_X \le 1} |\ell(v)|,
  \label{eq:ch02-dualnorm}
\end{equation}
and the evaluation of $\ell$ on $v$ is written $\langle \ell, v \rangle_{X^*, X}$, the \textbf{duality pairing}. The norm \eqref{eq:ch02-dualnorm} makes $X^*$ a Banach space whether or not $X$ is complete, since a Cauchy sequence of functionals converges pointwise, and it differs in kind from the norms met so far: it is defined by a supremum over a unit ball rather than by an integral, so it measures worst-case amplification rather than average size. Unlike an inner product, the pairing joins two \emph{different} spaces, and keeping them distinct is what lets a differential operator be inverted against rough test data. H\"older's inequality is the prototype: for $1 \le p, q \le \infty$ with $1/p + 1/q = 1$,
\begin{equation}
  \left| \int_\Omega u(x) v(x) \, dx \right| \le \|u\|_{L^p(\Omega)} \, \|v\|_{L^q(\Omega)},
  \label{eq:ch02-holder}
\end{equation}
so every $v \in L^q$ generates a functional on $L^p$ of norm at most $\|v\|_{L^q}$, and conversely, for $1 \le p < \infty$ every functional arises this way~\cite{brezis2011functional}:
\begin{equation}
  \big(L^p(\Omega)\big)^* \cong L^q(\Omega), \qquad \frac{1}{p} + \frac{1}{q} = 1, \quad 1 \le p < \infty .
  \label{eq:ch02-lpdual}
\end{equation}
The restriction $p < \infty$ in \eqref{eq:ch02-lpdual} is real and easy to lose track of. The dual of $L^\infty$ is strictly larger than $L^1$; it is the space of bounded finitely additive measures absolutely continuous with respect to Lebesgue measure, and $L^1$ sits inside it as a proper closed subspace~\cite{brezis2011functional}. Since the permeability coefficient of the Darcy problem in \chapref{ch:part1-partial-differential-equations-as-solution} is modeled in $L^\infty$, an $L^\infty$ bound on inputs does not by itself give any convergent subsequence in a norm topology: what it gives, through Banach--Alaoglu applied to $L^\infty = (L^1)^*$, is weak-$*$ compactness of the closed ball, a compactness that does not control the size of $\|a_n - a\|_{L^\infty}$ and therefore does not by itself deliver the uniform approximation property a network is trained to satisfy.

A Banach space is \textbf{reflexive} when the canonical embedding $J : X \to X^{**}$, $J(x)(\ell) = \ell(x)$, is onto. Hilbert spaces and, more generally, uniformly convex spaces are reflexive; so are $L^p(\Omega)$ and $W^{k,p}(\Omega)$ for $1 < p < \infty$~\cite{brezis2011functional}. Reflexivity is the hypothesis that makes ``bounded'' behave the way Heine--Borel makes it behave in finite dimensions, except that the compactness is weak rather than norm: the closed unit ball of a reflexive space is weakly compact. Existence proofs read this compactness off minimizing or Galerkin sequences, and any argument in this book that extracts a limit from a bound on $\|\nabla u\|_{L^2}$ is using it.

The construction that ties duality to regularity is the Gelfand triple, also called an evolution triple. One fixes three spaces with continuous injections, the middle one the pivot,
\begin{equation}
  V \hookrightarrow H \cong H^* \hookrightarrow V^* , \qquad
  V \text{ reflexive Banach}, \quad H \text{ Hilbert},
  \label{eq:ch02-triple}
\end{equation}
\index{evolution triple}\index{weak form!as a duality pairing}
where $V$ is dense in $H$, $H$ is identified with its own dual by Theorem~\ref{thm:ch02-riesz-rep}, and $V^*$ is the dual of $V$. The chain \eqref{eq:ch02-triple} must be read in one direction only: it is elements of $V$, and then of $H$, that act as functionals on $V$, and the identification $H \cong H^*$ is precisely what makes the restriction of $\langle \cdot, \cdot \rangle_{V^*, V}$ to $H \times V$ coincide with the inner product of $H$. A general element of $V^*$ has no value on a member of $H$ that lies outside $V$, and that asymmetry is the whole utility of the construction: the pairing extends the inner product downward in regularity rather than upward. The instance used throughout Part~\ref{part:1} is $V = H^1_0(\Omega)$, $H = L^2(\Omega)$, $V^* = H^{-1}(\Omega)$, and its importance is concrete: a source term $f$ that is not a function but a distribution still acts on $v \in H^1_0(\Omega)$ through this pairing, which is what permits the weak form of Darcy flow to be posed at all. Poincar\'e's inequality,
\begin{equation}
  \|v\|_{L^2(\Omega)} \le C_\Omega \|\nabla v\|_{L^2(\Omega)} \quad \forall v \in H^1_0(\Omega), \ \Omega \text{ bounded in one direction},
  \label{eq:ch02-poincare}
\end{equation}
supplies the equivalence of $\|\cdot\|_{H^1}$ with the Dirichlet seminorm on $H^1_0$ and is the coercivity that the Lax--Milgram argument of \chapref{ch:part1-partial-differential-equations-as-solution} invokes~\cite{evans2010partial}. Note what \eqref{eq:ch02-poincare} costs: it fails on $H^1(\Omega)$ functions with nonzero boundary trace, which is why pure-Neumann problems must be posed modulo constants and why the choice of space, not the choice of network, decides whether an elliptic problem is well posed.

% -----------------------------------------------------------------------------
\section{Sobolev Spaces and Compact Embeddings}
\index{weak derivative}\index{Sobolev space!of integer order}\index{Sobolev space!of fractional order}\index{Sobolev space!on the torus}

The derivative of a function that is not differentiable has to be defined against test functions rather than pointwise. For $\Omega \subset \mathbb{R}^d$ open and $u, v \in L^1_{\text{loc}}(\Omega)$, $v$ is the \textbf{$\alpha$-th weak derivative} of $u$, written $D^\alpha u = v$, if
\begin{equation}
  \int_\Omega u(x) D^\alpha \phi(x) \, dx = (-1)^{|\alpha|} \int_\Omega v(x) \phi(x) \, dx \quad \forall \phi \in C_c^\infty(\Omega).
  \label{eq:ch02-weakderiv}
\end{equation}
The identity \eqref{eq:ch02-weakderiv} is integration by parts with the derivative transferred onto the smooth compactly supported $\phi$, where it costs nothing; the definition's legitimacy rests on $C_c^\infty(\Omega)$ being dense enough that at most one $v \in L^1_{\mathrm{loc}}$ can satisfy it, so $D^\alpha u$ is unique almost everywhere. The example worth holding onto is $u(x) = |x|$ on $\mathbb{R}$, whose weak derivative is $\operatorname{sign}(x)$: kinks are admissible, jumps are not, and the distinction will reappear as a restriction on which Sobolev exponent a discontinuous solution can carry (Remark~\ref{rem:ch02-jumps}). The space is then
\begin{equation}
  W^{k,p}(\Omega) = \left\{ u \in L^p(\Omega) : D^\alpha u \in L^p(\Omega) \text{ for all } |\alpha| \le k \right\},
  \label{eq:ch02-wkp}
\end{equation}
with $k \in \mathbb{N}$ and $1 \le p \le \infty$; membership of \eqref{eq:ch02-wkp} asserts that every derivative up to order $k$, taken in the sense of \eqref{eq:ch02-weakderiv}, is again $p$-integrable. The space is made Banach by
\begin{equation}
  \|u\|_{W^{k,p}(\Omega)} = \left( \sum_{|\alpha| \le k} \|D^\alpha u\|_{L^p(\Omega)}^p \right)^{1/p},
  \label{eq:ch02-wkp-norm}
\end{equation}
the sum in \eqref{eq:ch02-wkp-norm} being replaced by the maximum over $|\alpha|$ when $p = \infty$. Because $W^{k,p}$ is a closed subspace of a finite product of complete spaces, it is complete; when $p = 2$ we write $H^k(\Omega) = W^{k,2}(\Omega)$, and the norm \eqref{eq:ch02-wkp-norm} is induced by
\begin{equation}
  \langle u, v \rangle_{H^k(\Omega)} = \sum_{|\alpha| \le k} \int_\Omega D^\alpha u(x) \, D^\alpha v(x) \, dx,
  \label{eq:ch02-hk-inner}
\end{equation}
so $H^k$ is Hilbert under the inner product \eqref{eq:ch02-hk-inner} and Theorem~\ref{thm:ch02-riesz-rep} applies to it, the representer reading out all derivatives up to order $k$ at once. The subspace that carries the boundary information is $H^1_0(\Omega)$, the closure of $C_c^\infty(\Omega)$ in $H^1(\Omega)$---roughly, the functions whose trace on $\partial\Omega$ vanishes---and it is the natural energy space for homogeneous Dirichlet problems.

For non-integer order a physical-space definition via difference quotients remains available, but on $\mathbb{T}^d$ it hides exactly the information one wants, and the Fourier characterization is cleaner and directly usable by the architectures of Part~\ref{part:2}. On the torus $\mathbb{T}^d = [0,1)^d$ any distribution $u$ possesses a Fourier expansion
\begin{equation}
  u(x) = \sum_{k \in \mathbb{Z}^d} \hat{u}(k) e^{2\pi i k \cdot x}, \qquad \hat{u}(k) = \int_{\mathbb{T}^d} u(x) e^{-2\pi i k \cdot x} \, dx,
  \label{eq:ch02-fourier-pair}
\end{equation}
and \eqref{eq:ch02-fourier-pair} converts differentiation into multiplication: $D^\alpha u$ has coefficients $(2\pi i k)^\alpha \hat{u}(k)$, so differentiability is nothing but polynomial decay of $\hat u$. Parseval's identity,
\begin{equation}
  \|u\|_{L^2(\mathbb{T}^d)}^2 = \sum_{k \in \mathbb{Z}^d} |\hat{u}(k)|^2,
  \label{eq:ch02-parseval}
\end{equation}
then motivates, for any $s \in \mathbb{R}$, the fractional Sobolev space $H^s(\mathbb{T}^d)$, whose elements are periodic distributions $u \in \mathcal{S}'(\mathbb{T}^d)$---the dual of $C^\infty(\mathbb{T}^d)$, the natural setting for a condition stated purely in terms of Fourier coefficients---with norm finite in the sense
\begin{equation}
  H^s(\mathbb{T}^d) = \left\{ u \in \mathcal{S}'(\mathbb{T}^d) : \|u\|_{H^s}^2 = \sum_{k \in \mathbb{Z}^d} (1 + |k|^2)^s |\hat{u}(k)|^2 < \infty \right\}.
  \label{eq:ch02-hs-torus}
\end{equation}
Definition \eqref{eq:ch02-hs-torus} is the one the rest of Part~\ref{part:1} uses, and three of its features do the work. First, $s$ may be negative: \eqref{eq:ch02-hs-torus} with $s = -1$ measures growth rather than decay of coefficients, and $H^{-s}(\mathbb{T}^d)$ is isometric to the dual of $H^s(\mathbb{T}^d)$ via the pairing
\begin{equation}
  \langle u, v \rangle_{H^{-s}, H^s} = \sum_{k \in \mathbb{Z}^d} \hat{u}(k) \, \overline{\hat{v}(k)} .
  \label{eq:ch02-hsneg}
\end{equation}
The duality in \eqref{eq:ch02-hsneg} extends Parseval's identity \eqref{eq:ch02-parseval} from square-integrable functions to distributions on the torus and is the exact analogue of the triple \eqref{eq:ch02-triple}: a residual small in $H^{-s}$ is small against every test field of unit $H^s$ mass, a far weaker demand than smallness in $L^2$. Second, the scale is interpolating: for $s < r < t$, writing $r = \theta s + (1-\theta)t$ with $0 < \theta < 1$, the pointwise bound on the weights combined with H\"older's inequality \eqref{eq:ch02-holder} applied to counting measure on $\mathbb{Z}^d$ gives
\begin{equation}
  \|u\|_{H^r} \le \|u\|_{H^s}^{\theta} \, \|u\|_{H^t}^{1-\theta},
  \label{eq:ch02-interp}
\end{equation}
so a bound at two levels of regularity controls every level between them; interpolation \eqref{eq:ch02-interp} is the step taken whenever an estimate must be degraded from the strong norm in which it was proved to the weaker norm a loss function actually sees. Third, and most concretely, \eqref{eq:ch02-hs-torus} makes pointwise evaluation decidable. For $s > d/2$ the series in \eqref{eq:ch02-fourier-pair} converges absolutely and uniformly, so each $u \in H^s(\mathbb{T}^d)$ has a continuous representative and its point values are well defined; the triangle inequality and Cauchy--Schwarz then yield
\begin{equation}
  |u(x)| \le \sum_{k \in \mathbb{Z}^d} |\hat{u}(k)|
  = \sum_{k} (1+|k|^2)^{-s/2} \, (1+|k|^2)^{s/2} |\hat{u}(k)|
  \le C_{s,d} \, \|u\|_{H^s(\mathbb{T}^d)},
  \label{eq:ch02-eval}
\end{equation}
with $C_{s,d}^2 = \sum_k (1+|k|^2)^{-s} < \infty$ precisely when $2s > d$. The bound \eqref{eq:ch02-eval} is the honest reason Deep Operator Networks (DeepONet) and other sensor-based architectures can be given a theory at all: the branch network consumes samples $u(x_j)$, and those are continuous linear functionals only on spaces carrying more than $d/2$ derivatives in the $L^2$ sense. On $L^2(\mathbb{T}^d)$ itself a point evaluation is unbounded, so a branch network reading nodal values has no well-defined continuum limit; the sensor term of the DeepONet error analysis of Lanthaler et al.~\cite{lanthaler2022error}, priced against a finite number $m$ of readings in \chapref{ch:part1-operator-approximation-theory-and-universa}, counts the cost of that finite budget, while \eqref{eq:ch02-eval} decides whether the readings are bounded functionals at all. The threshold $s > d/2$ reappears in \eqref{eq:ch02-linfty-embed} below and in the loss design of \chapref{ch:part1-physics-informed-operators-and-sobolev-reg}; it is worth internalizing once, carefully, rather than three times approximately.
\index{DeepONet (Deep Operator Network)!pointwise sensor evaluations}\index{point evaluation functional}

% -----------------------------------------------------------------------------
\section{Embedding Theorems and the Cost of Compactness}

The embeddings collect these comparisons between spaces, with their hypotheses exposed, because the convergence results of \chapref{ch:part1-the-fourier-neural-operator-spectral-theor} and \chapref{ch:part1-physics-informed-operators-and-sobolev-reg} use them silently.

\begin{proposition}[Sobolev embeddings, explicit conditions]
\label{prop:ch02-sobolev-embeddings}
Let $\Omega \subset \mathbb{R}^d$ be bounded with $C^1$ boundary, let $k \in \mathbb{N}_0$ and $1 \le p < \infty$. In the subcritical regime $kp < d$ one has $W^{k,p}(\Omega) \hookrightarrow L^q(\Omega)$ continuously for every $q$ with
\begin{equation}
  1 \le q \le p^{*}, \qquad p^{*} := \frac{dp}{d - kp}, \qquad \text{with} \quad \frac{1}{q} = \frac{1}{p} - \frac{k}{d} \ \text{at} \ q = p^{*},
  \label{eq:ch02-sobembed}
\end{equation}
and the embedding is compact whenever $q < p^{*}$. At the borderline $kp = d$ the target range widens to every finite $q$, that is $W^{k,p}(\Omega) \hookrightarrow L^q(\Omega)$ continuously and compactly for all $1 \le q < \infty$, but no embedding into $L^\infty(\Omega)$ holds. In the supercritical regime $kp > d$, writing $\sigma := k - d/p$, one has $W^{k,p}(\Omega) \hookrightarrow C^{j,\alpha}(\overline{\Omega})$ continuously for every integer $j \ge 0$ and $0 < \alpha < 1$ with $j + \alpha \le \sigma$, the endpoint $j + \alpha = \sigma$ being admissible only when $\sigma$ is not an integer; the embedding is compact whenever $j + \alpha$ is strictly below $\sigma$. On $\mathbb{T}^d$ the same three alternatives hold with $k$ replaced by a fractional order $s$, and in particular
\begin{equation}
  H^s(\mathbb{T}^d) \hookrightarrow L^\infty(\mathbb{T}^d) \quad \text{continuously if and only if} \quad s > \frac{d}{2}.
  \label{eq:ch02-linfty-embed}
\end{equation}
\end{proposition}

The proof is a transplant argument: the hypothesis on $\partial\Omega$ provides an extension operator $W^{k,p}(\Omega) \to W^{k,p}(\mathbb{R}^d)$, after which \eqref{eq:ch02-sobembed} follows from convolution estimates, with \eqref{eq:ch02-eval} giving the Hilbert case directly; both developments appear in Brezis~\cite{brezis2011functional} and Evans~\cite{evans2010partial}. Here $C^{j,\alpha}(\overline{\Omega})$ denotes the space of $j$ times continuously differentiable functions whose $j$-th derivatives are $\alpha$-H\"older continuous. The $C^1$ assumption on $\partial\Omega$ is more than these proofs need: the extension step holds on any domain with the strong local Lipschitz property, which is the class the notation table admits, and no argument in this Part requires the refinement. The supercritical clause is the only route by which an integral bound can imply a pointwise one, and it is the mechanism behind the $L^\infty$ convergence claim attached to Sobolev-regularized losses later in this Part. The condition \eqref{eq:ch02-linfty-embed} is also sharp in exactly the way the borderline case warns: a bounded set in $H^{d/2}(\mathbb{T}^d)$ need not be bounded in $L^\infty$, so a loss that controls precisely $d/2$ derivatives buys nothing pointwise. Insisting on $s > d/2$ is not fastidiousness.
\index{Sobolev embedding theorem}\index{embedding!continuous}\index{embedding!compact}\index{Morrey-type embedding}

Compactness of the injection is the second half of the story, and it is the theorem on which the approximation theory of the next two chapters rests.

\begin{theorem}[Rellich--Kondrachov Compact Embedding]
\label{thm:ch02-rellich}
Let $\Omega \subset \mathbb{R}^d$ be bounded with $C^1$ boundary. If $s > t \ge 0$, the canonical injection
\begin{equation}
  \iota : H^s(\Omega) \hookrightarrow H^t(\Omega)
  \label{eq:ch02-injection}
\end{equation}
is a compact linear operator. Furthermore, if $s > \frac{d}{2} + k$, then $H^s(\Omega) \hookrightarrow C^k(\overline{\Omega})$ continuously (Sobolev Embedding Theorem).
\end{theorem}

On the torus the content of \eqref{eq:ch02-injection} is transparent and worth seeing, because it is the same computation that produces the spectral convergence rate in \chapref{ch:part1-the-fourier-neural-operator-spectral-theor}. Let $P_K$ denote the orthogonal projection of $L^2(\mathbb{T}^d)$ onto the span of the modes with $|k_j| \le K$ for every $j$, and let $w \in H^\sigma(\mathbb{T}^d)$ with $\sigma > t \ge 0$. Since $(1+|k|^2)^{t} \le (1+K^2)^{t-\sigma}(1+|k|^2)^{\sigma}$ on the discarded modes $|k|_\infty > K$, where $t - \sigma < 0$,
\begin{equation}
  \|w - P_K w\|_{H^t(\mathbb{T}^d)}^2
  = \sum_{|k|_\infty > K} (1+|k|^2)^{t} |\hat{w}(k)|^2
  \le (1+K^2)^{t-\sigma} \, \|w\|_{H^\sigma(\mathbb{T}^d)}^2 .
  \label{eq:ch02-tail}
\end{equation}
Now take a set bounded in $H^s(\mathbb{T}^d)$, $s > t$, say by $M$, and fix $\varepsilon > 0$. Choosing $K$ with $(1+K^2)^{(t-s)/2} M \le \varepsilon$, \eqref{eq:ch02-tail} puts every element of the image within $\varepsilon$ of its own projection $P_K u$, and those projections form a bounded set in the finite-dimensional space $\operatorname{span}\{e^{2\pi i k \cdot x} : |k|_\infty \le K\}$, hence a totally bounded one. Total boundedness survives $\varepsilon$-approximation, so the image is totally bounded in $H^t$, and completeness of $H^t$ upgrades that to relative compactness---which is \eqref{eq:ch02-injection} in the periodic case, obtained with no boundary regularity at all. The estimate \eqref{eq:ch02-tail} is the quantitative residue of the argument: the rate $(1+K^2)^{(t-\sigma)/2}$ is the algebraic truncation decay that the mode-cutoff error of the Fourier neural operator (FNO) inherits in \chapref{ch:part1-the-fourier-neural-operator-spectral-theor}~\cite{kovachki2021universal}.
\index{Rellich--Kondrachov theorem}\index{compact embedding!proof on the torus}\index{spectral truncation!tail estimate}

The relevance to learning is immediate, and it deserves to be stated precisely. If an input ensemble, or a PDE solution operator, maps bounded sets into a space of higher Sobolev regularity---$H^s \to H^{s+\alpha}$ with $\alpha > 0$---then its image on a bounded set is precompact in $H^s$ by Theorem~\ref{thm:ch02-rellich}, and a uniform approximation theorem whose domain must be compact can be applied~\cite{kovachki2023neural}. Parabolic and elliptic smoothing supply such $\alpha$ generously, as \chapref{ch:part1-partial-differential-equations-as-solution} shows case by case; hyperbolic propagation does not. Two limits deserve recording now rather than discovery later. \eqref{eq:ch02-injection} is qualitative: it certifies that uniform approximation is possible, not how cheaply, and the sample and width complexity of the resulting network is a separate question that Part~\ref{part:1} answers only for specific architectures. And the smoothing order $\alpha$ is a hypothesis about the operator being learned, not a property of the architecture, so a model validated on a smoothing problem carries no guarantee on a transport problem with discontinuous data.

\begin{remark}[Why discontinuities break the scale]
\label{rem:ch02-jumps}
In one space dimension a function with a jump discontinuity has Fourier coefficients of size $|\hat{u}(k)| \sim |k|^{-1}$, so by \eqref{eq:ch02-hs-torus} it belongs to $H^s(\mathbb{T})$ for $s < \tfrac{1}{2}$ and to no $H^s$ with $s \ge \tfrac{1}{2}$. A family of solutions whose jumps and $L^2$ norms are uniformly bounded is therefore bounded in $H^s(\mathbb{T})$ for $s < \tfrac{1}{2}$, and generically unbounded in every $H^s$ with $s \ge \tfrac{1}{2}$; Theorem~\ref{thm:ch02-rellich} still yields compactness in $H^t$ for $t < s$, but the usable $t$ sits well below the regularity a spectral method wants. What fails is the hypothesis of the Kolmogorov--Riesz criterion quoted earlier: a bound in $W^{1,1}$ controls translations through $\|u(\cdot+h)-u\|_{L^1} \le |h| \|\nabla u\|_{L^1}$, and no such control survives an $L^\infty$ bound alone. This is the functional-analytic content of the Gibbs-phenomenon and shock-capturing difficulties recorded in \chapref{ch:part1-open-theoretical-problems-and-mathematical}.
\end{remark}

\takeaway{What This Chapter Is For}{Every error statement in this book is a norm, every norm comes with a compactness criterion, and every compactness criterion in infinite dimensions costs regularity. Sobolev exponent $s$ is the currency: $s > d/2$ makes pointwise values and $L^\infty$ errors controllable, $s > t$ makes $H^s$-bounded sets precompact in $H^t$, and a smoothing gain of $\alpha > 0$ derivatives is what converts a bounded input set into a compact one that a uniform approximation theorem can act on.}

% -----------------------------------------------------------------------------
\section{Bochner Spaces for Time-Dependent Evolution}
\index{Bochner space}\index{strong measurability}\index{Aubin--Lions--Simon lemma}

Evolution problems ask for a function of time whose values are functions of space, and the right object is a single point in a function space over a function space. For a Banach space $X$ and $1 \le p \le \infty$, the space $L^p(0,T;X)$ consists of (equivalence classes of) strongly measurable maps $u : [0,T] \to X$---limits almost everywhere of simple functions, which for separable $X$ coincides with ordinary measurability of $u : [0,T] \to X$~\cite{brezis2011functional}---with norm
\begin{equation}
  \|u\|_{L^p(0, T; X)} = \left( \int_0^T \|u(t)\|_X^p \, dt \right)^{1/p} < \infty,
  \label{eq:ch02-bochner}
\end{equation}
the usual essential-supremum modification when $p = \infty$. The definition \eqref{eq:ch02-bochner} does two things at once. It turns a trajectory $t \mapsto u(\cdot,t)$ into one element of one Banach space, so that a solution operator for an evolution equation can be treated as a bounded nonlinear map between \emph{static} spaces and the machinery of the preceding sections applies to it unchanged; and it makes the time variable structurally identical to the space variable, which is why a surrogate that learns a whole trajectory and one that learns a single time slice are analyzed by the same theorems but not at the same rates. Completeness holds exactly when $X$ is complete, and duality is as good as one could hope: for $1 < p < \infty$ with $X$ reflexive,
\begin{equation}
  \big(L^p(0,T;X)\big)^* \cong L^q(0,T;X^*), \qquad
  \langle\!\langle \ell, u \rangle\!\rangle = \int_0^T \langle \ell(t), u(t) \rangle_{X^*, X} \, dt ,
  \label{eq:ch02-bochner-dual}
\end{equation}
so the representation \eqref{eq:ch02-bochner-dual} is what lets the weak form of an evolutionary PDE be tested against space-time functions in $L^2(0,T;H^1_0(\Omega))$, whose dual is $L^2(0,T;H^{-1}(\Omega))$~\cite{evans2010partial}. \index{duality!of Bochner spaces}

Weak time derivatives are defined exactly as in \eqref{eq:ch02-weakderiv}, with the test functions taken in $C_c^\infty(0,T)$ and values in $X$, and the natural energy space for a second-order parabolic problem on the triple \eqref{eq:ch02-triple} is the intersection
\begin{equation}
  \begin{aligned}
    W(0,T) &:= L^2\big(0,T; H^1_0(\Omega)\big) \cap H^1\big(0,T; H^{-1}(\Omega)\big), \\
    \|u\|_{W(0,T)} &= \|u\|_{L^2(0,T;H^1_0)} + \|\partial_t u\|_{L^2(0,T;H^{-1})}.
  \end{aligned}
  \label{eq:ch02-w0t}
\end{equation}
The norm \eqref{eq:ch02-w0t} makes $W(0,T)$ a Hilbert space, and its structure is the point: $u$ itself has one spatial derivative in $L^2$, while $\partial_t u$ is allowed to be a distribution in space, measured in $H^{-1}$. That asymmetry is not a defect but an accurate description of what diffusion estimates deliver, and it is also the reason continuity in time cannot simply be read off---one needs a compactness lemma. The Aubin--Lions--Simon lemma supplies one: whenever a Banach space $X$ embeds compactly into a reflexive space $B$ and $B$ embeds continuously into a Banach space $F$, the set of uniformly bounded trajectories
\begin{equation}
  \left\{ u \in L^p(0,T;X) : \|u\|_{L^p(0,T;X)} \le C, \ \|\partial_t u\|_{L^r(0,T;F)} \le C \right\}
  \quad \text{is relatively compact in } L^p(0,T;B),
  \label{eq:ch02-aubinlions}
\end{equation}
for $1 \le p < \infty$ and $1 \le r \le \infty$, and, in the limiting case $p = \infty$ with $r > 1$, relatively compact in $C([0,T];B)$~\cite{evans2010partial}. Taking $X = H^1_0(\Omega)$, $B = L^2(\Omega)$, $F = H^{-1}(\Omega)$, where the first injection is Theorem~\ref{thm:ch02-rellich} and the second is the dense continuous embedding produced by the triple \eqref{eq:ch02-triple}, bounded subsets of the space defined in \eqref{eq:ch02-w0t} are relatively compact in $L^2(0,T;L^2(\Omega))$; the same triple also gives $u \in C([0,T];L^2(\Omega))$ after redefinition on a null set of times, together with the energy identity
\begin{equation}
  \frac{1}{2}\frac{d}{dt}\|u(t)\|_{L^2(\Omega)}^2 = \big\langle \partial_t u(t), u(t) \big\rangle_{H^{-1}, H^1_0}
  \quad \text{for a.e. } t \in (0,T),
  \label{eq:ch02-energy}
\end{equation}
which is the statement that testing the equation against the solution itself is legitimate~\cite{evans2010partial}. The identity \eqref{eq:ch02-energy} is the entry point for every dissipation estimate in this book, including the ones that justify the smoothing hypothesis on which the application of Theorem~\ref{thm:ch02-rellich} above rests.

For operator learning the payoff is a change of viewpoint with consequences. The map $u_0 \mapsto u(\cdot,\cdot)$ that sends initial data to a space-time solution is, for a smoothing parabolic problem, a bounded map $H^s(\Omega) \to L^2(0,T;H^{s+\alpha}(\Omega))$. That bound by itself buys no compactness in time: with $\varphi \in H^{s+\alpha}(\Omega)$ fixed, the sequence $\sin(nt)\,\varphi$ is bounded in every $L^p(0,T;H^{s+\alpha}(\Omega))$ and has no convergent subsequence in $L^2(0,T;H^s(\Omega))$. The hypothesis of \eqref{eq:ch02-aubinlions} it fails to check is the one on $\partial_t u$, and the equation, not the architecture, supplies it: for a second-order parabolic problem the energy estimate that bounds $u$ in $L^2(0,T;H^{s+1}(\Omega))$ bounds $\partial_t u$ in $L^2(0,T;H^{s-1}(\Omega))$ as well, so the triple $H^{s+1}(\Omega) \hookrightarrow H^{s}(\Omega) \hookrightarrow H^{s-1}(\Omega)$ is admitted by \eqref{eq:ch02-aubinlions} with $\alpha = 1$, and the image of a bounded set of initial data is precompact in $L^2(0,T;H^s(\Omega))$---Theorem~\ref{thm:ch02-rellich} in space, \eqref{eq:ch02-aubinlions} in time. That is the setting in which an autoregressive stepping operator $u(t) \mapsto u(t+\Delta t)$, or a trajectory-to-trajectory model, can be discussed at all. What the framework does \emph{not} give is stability of the learned recursion: \eqref{eq:ch02-aubinlions} and \eqref{eq:ch02-energy} constrain the true solution operator, and nothing constrains a surrogate to inherit the dissipation that makes \eqref{eq:ch02-energy} integrable over long times. The mismatch between a compactness statement about the target and a long-horizon rollout of an approximation is left open here and is taken up in \chapref{ch:part1-open-theoretical-problems-and-mathematical}. \index{autoregressive rollout}

% -----------------------------------------------------------------------------
This chapter fixed the spaces, the norms, and the compactness certificates, and it was careful to charge every later argument its hypotheses: a smoothing gain $\alpha > 0$, a regularity threshold $s > d/2$ for anything pointwise, boundary regularity on $\Omega$ for the extension step behind Proposition~\ref{prop:ch02-sobolev-embeddings} and Theorem~\ref{thm:ch02-rellich}, and reflexivity or the range $1 < p < \infty$ whenever a dual was identified with an $L^q$ space. It also left the central question untouched. We know how to measure $\|\mathcal{G}(a_1) - \mathcal{G}(a_2)\|_{\mathcal{U}}$ and we know which subsets of $\mathcal{A}$ are compact enough for a uniform approximation theorem to bite, but nothing so far establishes that a specific physical law gives a map that is bounded, continuous, or differentiable between those spaces, let alone that it smooths by a definite order $\alpha$. That is the substance of the next chapter, \chapref{ch:part1-partial-differential-equations-as-solution}, which turns the Darcy, Burgers, and Navier--Stokes problems from equations into operators. There the weak forms will be written as pairings in the triple \eqref{eq:ch02-triple}, existence will rest on the coercivity supplied by \eqref{eq:ch02-poincare}, and the Lipschitz and Fr\'echet properties of the solution map will be checked against exactly the norms introduced here---so that when an architecture is later asked to learn one of these operators, we will know which of the assumptions made in this chapter it is being asked to inherit.
